\section{Conclusion}
\label{sec:conclusion}
In this paper, we introduce \ourmethod, a branch-and-merge based multi-agent system for long-horizon software engineering tasks. We use a manager to break a task into dependency-aware units, assign them to engineers, and keep each engineer working in an isolated branch and worktree. Progress is integrated only through \texttt{git commit} and \texttt{git merge} on the main branch, with tests used as the executable check for whether an update should be kept. Across Commit0 and PaperBench, our \ourmethod consistently improves over single-agent baselines, even when the underlying model is unchanged. Our results also show that simply increasing the single agent iteration budget does not reliably improve outcomes, and a fallback strategy that runs a single agent first and then switches to multi-agent mainly wastes runtime and cost. Overall, we show that branch-and-merge is important for effective multi-agent software engineering and that SWE primitives provide a practical way to support it. For complex long-horizon, dependency-aware software engineering tasks, \ourmethod is the default paradigm for structuring solutions to enable parallel and coordinated development.

% \gncomment{We might want a slightly stronger statement here. For instance, something that indicates ``if you are doing XXX type of task, you should be using our method''.}

% Overall,  for long-horizon tasks with shared artifacts and integration gates, performance depends heavily on how work is organized, delegated and merged into a stable repository state, and branch-and-merge based coordination provides an effective way to do that.